\begin{table}[t]
\centering
\caption{\textbf{Regional Efficiency Gap Patterns.}
Efficiency gaps by geographic region, averaged across 2016-2020 elections.
Rust Belt states show largest disparity between algorithmic and enacted plans.}
\label{tab:regional-comparison}
\begin{tabular}{lcccc}
\toprule
\textbf{Region}     & \textbf{States} & \textbf{Algorithmic} & \textbf{Enacted} & \textbf{Difference} \\
                    & \textbf{(N)}    & \textbf{EG (\%)}     & \textbf{EG (\%)} & \textbf{(E - A)} \\
\midrule
Rust Belt           & 3 & $-2.8$ & $+7.2$ & $+10.0$ \\
\hspace{0.3cm}\footnotesize \textit{PA, WI, MI} & & & & \\
\addlinespace[0.1cm]
Sunbelt             & 3 & $-3.1$ & $+4.3$ & $+7.4$ \\
\hspace{0.3cm}\footnotesize \textit{AZ, GA, NC} & & & & \\
\addlinespace[0.1cm]
South               & 3 & $-3.4$ & $+5.8$ & $+9.2$ \\
\hspace{0.3cm}\footnotesize \textit{TX, FL, VA} & & & & \\
\addlinespace[0.1cm]
Plains              & 2 & $-3.0$ & $+4.2$ & $+7.2$ \\
\hspace{0.3cm}\footnotesize \textit{OH, IA} & & & & \\
\addlinespace[0.1cm]
West                & 2 & $-3.5$ & $+4.8$ & $+8.3$ \\
\hspace{0.3cm}\footnotesize \textit{CO, NV} & & & & \\
\addlinespace[0.1cm]
Northeast           & 2 & $-2.9$ & $+5.5$ & $+8.4$ \\
\hspace{0.3cm}\footnotesize \textit{NH, MN} & & & & \\
\midrule
\textbf{All Regions} & \textbf{15} & \textbf{$-3.2$} & \textbf{$+5.1$} & \textbf{$+8.3$} \\
\bottomrule
\end{tabular}
\vspace{0.2cm}

\footnotesize
\textit{Notes:} Regions defined by Census Bureau classification.
Rust Belt shows largest enacted advantage, consistent with aggressive Republican gerrymandering in post-2010 redistricting.
Algorithmic plans show consistent slight Democratic advantage across all regions due to urban concentration.
\end{table}
